\documentclass[a4paper,12pt,titlepage,finall]{article}

\usepackage[T1,T2A]{fontenc}
\usepackage[utf8x]{inputenc}
\usepackage[russian]{babel}
\usepackage{tikz}
\usepackage{pgfplots}
\usepackage{geometry}
\usepackage{indentfirst}
\usepackage{amsmath,amssymb}
\usepackage{array}
\usepackage{longtable}
\usepackage{url}
\usepackage{float}

\geometry{a4paper,left=30mm,top=30mm,bottom=30mm,right=30mm}
\setcounter{secnumdepth}{0}
\usepgfplotslibrary{fillbetween}
\pgfplotsset{compat=1.18}

\begin{document}

\begin{titlepage}
    \begin{center}
    {\small \sc Московский государственный университет \\имени М.~В.~Ломоносова\\
    Факультет вычислительной математики и кибернетики\\}
    \vfill
    {\Large \sc Отчет по заданию №6}\\
    ~\\
    {\large \bf <<Сборка многомодульных программ. \\
    Вычисление корней уравнений и определенных интегралов>>}\\
    ~\\
    {\large \bf Вариант 2 / 2,3 / 2}
    \end{center}
    \begin{flushright}
    \vfill {Выполнила:\\
    студентка 104 группы\\
    Грознецких~Екатерина~Дмитриевна\\
    ~\\
    Преподаватель:\\
    Гуляев~Денис~Анатольевич\\
    }
    \end{flushright}
    \begin{center}
    \vfill
    {\small Москва\\2026}
    \end{center}
\end{titlepage}

\tableofcontents

\newpage
\section{Постановка задачи}

Требуется реализовать многомодульную программу на языке Си и языке ассемблера NASM32, вычисляющую площадь плоской фигуры, ограниченной графиками трех функций. В соответствии с вариантом 2 заданы функции
\[
    f_1(x)=3\left(\frac{0.5}{x+1}+1\right),\qquad
    f_2(x)=2.5x-9.5,\qquad
    f_3(x)=\frac{5}{x},\quad x>0.
\]
В программе используется тип данных \texttt{long double}. Функции \(f_1\), \(f_2\), \(f_3\), а также их производные реализованы в отдельном ассемблерном модуле \texttt{funcs.asm}. Остальная часть программы, включая поиск корней, вычисление интегралов, обработку ключей командной строки и вывод результатов, реализована в модуле \texttt{main.c}.

Вершины фигуры находятся как точки пересечения пар кривых. Для этого решаются уравнения
\[
    f_1(x)=f_3(x),\qquad f_2(x)=f_3(x),\qquad f_1(x)=f_2(x).
\]
Поиск корней выполняется комбинированным методом: на каждом шаге используется метод хорд и метод Ньютона. Для вычисления площади применяется метод трапеций. В экспериментах используются значения точности
\[
    \varepsilon=10^{-3},\qquad \varepsilon=10^{-5},\qquad \varepsilon=10^{-7}.
\]

Отрезки, на которых применяются численные методы нахождения корней, предварительно выбираются аналитически по виду уравнений и знакам разностей функций на концах отрезков. В программе использованы отрезки \([1,2]\), \([4,5]\) и \([5,6]\).

\newpage

\section{Математическое обоснование}

Заданная первая функция может быть записана в более удобном виде:
\[
    f_1(x)=3\left(\frac{0.5}{x+1}+1\right)=3+\frac{1.5}{x+1}.
\]
Производные функций, необходимые для метода Ньютона, равны
\[
    f'_1(x)=-\frac{1.5}{(x+1)^2},\qquad
    f'_2(x)=2.5,
\]
\[
    f'_3(x)=-\frac{5}{x^2}.
\]
Все эти функции определены на рассматриваемых отрезках, причем для \(f_3\) дополнительно учитывается ограничение \(x>0\).

Для выбора отрезков поиска точек пересечения были аналитически рассмотрены уравнения пересечения пар функций.

Для пересечения \(f_1\) и \(f_3\):
\[
    3+\frac{1.5}{x+1}=\frac{5}{x}.
\]
После умножения на \(x(x+1)\) получаем
\[
    3x^2-0.5x-5=0,
\]
или
\[
    6x^2-x-10=0.
\]
Положительный корень равен
\[
    x_{13}=\frac{1+\sqrt{241}}{12}\approx 1.3770.
\]
Поэтому для поиска этого корня выбран отрезок \([1,2]\).

Для пересечения \(f_2\) и \(f_3\):
\[
    2.5x-9.5=\frac{5}{x}.
\]
После умножения на \(x\) получаем
\[
    2.5x^2-9.5x-5=0,
\]
или
\[
    5x^2-19x-10=0.
\]
Положительный корень равен
\[
    x_{23}=\frac{19+\sqrt{561}}{10}\approx 4.2685.
\]
Поэтому выбран отрезок \([4,5]\).

Для пересечения \(f_1\) и \(f_2\):
\[
    3+\frac{1.5}{x+1}=2.5x-9.5.
\]
Отсюда
\[
    (2.5x-12.5)(x+1)-1.5=0,
\]
то есть
\[
    2.5x^2-10x-14=0,
\]
или
\[
    5x^2-20x-28=0.
\]
Положительный корень равен
\[
    x_{12}=2+0.8\sqrt{15}\approx 5.0984.
\]
Поэтому выбран отрезок \([5,6]\).

Графики функций приведены на рис.~\ref{plot1}. На рассматриваемой фигуре верхней границей является функция \(f_1\). На отрезке \([x_{13},x_{23}]\) нижней границей является \(f_3\), а на отрезке \([x_{23},x_{12}]\) нижней границей является \(f_2\). Поэтому площадь фигуры вычисляется по формуле
\[
    S=\int_{x_{13}}^{x_{23}}\left(f_1(x)-f_3(x)\right)\,dx+
      \int_{x_{23}}^{x_{12}}\left(f_1(x)-f_2(x)\right)\,dx.
\]

Для метода Ньютона требуется, чтобы в окрестности корня функция была дифференцируема, а производная не обращалась в ноль. В программе метод применяется к разности двух функций, то есть к функции \(F(x)=f(x)-g(x)\). Производная этой функции равна \(F'(x)=f'(x)-g'(x)\). На выбранных отрезках функции непрерывны, их производные существуют, а сами отрезки содержат по одному корню. Метод хорд дополнительно использует значения функции на концах отрезка, где разность функций меняет знак.

В реализации одно значение \(\varepsilon\) используется и для поиска корней, и для остановки уточнения интеграла. Для метода трапеций число разбиений удваивается до тех пор, пока модуль разности двух последовательных приближений интеграла не станет меньше \(\varepsilon\).

\begin{figure}[H]
\centering
\begin{tikzpicture}
\begin{axis}[
             grid=both,
             axis lines=middle,
             restrict x to domain=0.5:5.6,
             restrict y to domain=-1:5,
             enlargelimits,
             legend cell align=left,
             width=13cm,
             height=9cm,
             xtick={0.5,1,2,3,4,5,6},
             xticklabel style={font=\small}]

\addplot[green,samples=256,thick,domain=0.5:5.6] {3 + 1.5/(x+1)};
\addlegendentry{$f_1(x)=3+\frac{1.5}{x+1}$}

\addplot[blue,samples=256,thick,domain=0.5:5.6] {2.5*x - 9.5};
\addlegendentry{$f_2(x)=2.5x-9.5$}

\addplot[red,samples=256,thick,domain=0.5:5.6] {5/x};
\addlegendentry{$f_3(x)=\frac{5}{x}$}
\end{axis}
\end{tikzpicture}
\caption{Графики функций варианта 2}
\label{plot1}
\end{figure}


\clearpage

\section{Результаты экспериментов}

В результате работы программы были найдены точки пересечения кривых и площадь ограниченной ими фигуры. Координаты точек пересечения приведены в таблице~\ref{table1}. Значения \(y\) вычислены подстановкой найденных абсцисс в соответствующие функции.

\begin{table}[h]
\centering
\begin{tabular}{|c|c|c|}
\hline
Кривые & $x$ & $y$ \\
\hline
1 и 3 & 1.377015 & 3.631044 \\
2 и 3 & 4.268544 & 1.171360 \\
1 и 2 & 5.098387 & 3.245967 \\
\hline
\end{tabular}
\caption{Координаты точек пересечения}
\label{table1}
\end{table}

Результаты вычислений для различных значений точности приведены в таблице~\ref{table2}. В таблице также указано количество итераций, потребовавшееся функции \texttt{root} для нахождения каждой точки пересечения.

\begin{table}[h]
\centering
\begin{tabular}{|c|c|c|c|c|c|c|c|}
\hline
$\varepsilon$ & $x_{13}$ & ит. & $x_{23}$ & ит. & $x_{12}$ & ит. & $S$ \\
\hline
$10^{-3}$ & 1.377182 & 6 & 4.268557 & 2 & 5.098387 & 2 & 5.087822 \\
$10^{-5}$ & 1.377018 & 9 & 4.268544 & 3 & 5.098387 & 2 & 5.087770 \\
$10^{-7}$ & 1.377015 & 13 & 4.268544 & 4 & 5.098387 & 2 & 5.087767 \\
\hline
\end{tabular}
\caption{Результаты работы программы}
\label{table2}
\end{table}

При уменьшении \(\varepsilon\) значения корней и площади стабилизируются. Итоговое значение площади при \(\varepsilon=10^{-7}\):
\[
    S\approx 5.087767.
\]

На рис.~\ref{plot2} показана ограниченная область, площадь которой вычислялась программой.

\begin{figure}[H]
\centering
\begin{tikzpicture}
\begin{axis}[
             grid=both,
             axis lines=middle,
             restrict x to domain=0.5:5.6,
             restrict y to domain=-1:5,
             enlargelimits,
             legend cell align=left,
             width=13cm,
             height=9cm,
             xtick={0.5,1,2,3,4,5,6},
             xticklabel style={font=\small}]

\addplot[green,samples=256,thick,domain=0.5:5.6,name path=A] {3 + 1.5/(x+1)};
\addlegendentry{$f_1(x)$}

\addplot[blue,samples=256,thick,domain=0.5:5.6,name path=B] {2.5*x - 9.5};
\addlegendentry{$f_2(x)$}

\addplot[red,samples=256,thick,domain=0.5:5.6,name path=C] {5/x};
\addlegendentry{$f_3(x)$}

\addplot[blue!20,samples=256] fill between[of=A and C,soft clip={domain=1.377015:4.268544}];
\addplot[blue!20,samples=256] fill between[of=A and B,soft clip={domain=4.268544:5.098387}];
\addlegendentry{$S\approx 5.087767$}

\addplot[dashed] coordinates {(1.377015,3.631044) (1.377015,0)};
\node[font=\scriptsize, fill=white, inner sep=1pt, anchor=south] at (axis cs:1.377015,0.15) {$x_{13}=1.377015$};

\addplot[dashed] coordinates {(4.268544,1.171360) (4.268544,0)};
\node[font=\scriptsize, fill=white, inner sep=1pt, anchor=south west] at (axis cs:4.268544,0.15) {$x_{23}=4.268544$};

\addplot[dashed] coordinates {(5.098387,3.245967) (5.098387,0)};
\node[font=\scriptsize, fill=white, inner sep=1pt, anchor=south west] at (axis cs:5.098387,0.45) {$x_{12}=5.098387$};

\end{axis}
\end{tikzpicture}
\caption{Плоская фигура, ограниченная графиками заданных функций}
\label{plot2}
\end{figure}

\clearpage

\section{Структура программы и спецификация функций}

Программа состоит из трех файлов:
\begin{itemize}
    \item \texttt{main.c} --- основной модуль на языке Си;
    \item \texttt{funcs.asm} --- модуль на ассемблере NASM32 с функциями и производными;
    \item \texttt{Makefile} --- файл сборки проекта.
\end{itemize}

Общая схема взаимодействия модулей показана ниже.
\[
\begin{array}{c}
\boxed{\texttt{main.c}}\\[2mm]
\downarrow\ \text{вызывает функции и производные}\\[2mm]
\boxed{\texttt{funcs.asm}}\\[2mm]
\end{array}
\qquad
\begin{array}{c}
\boxed{\texttt{Makefile}}\\[2mm]
\downarrow\\[2mm]
\text{компиляция } \texttt{main.c} \text{ и } \texttt{funcs.asm}
\end{array}
\]

Ниже приведена спецификация функций программы.

\begin{longtable}{|p{0.30\textwidth}|p{0.62\textwidth}|}
\hline
Функция & Назначение \\
\hline
\texttt{f1(long double x)} & Ассемблерная функция. Возвращает значение \(f_1(x)=3+\frac{1.5}{x+1}\). Значение аргумента \(x\) передается из Си-кода, результат возвращается через регистр \(st(0)\) сопроцессора x87. \\
\hline
\texttt{f2(long double x)} & Ассемблерная функция. Возвращает значение \(f_2(x)=2.5x-9.5\). Используется при поиске корней, вычислении интегралов и тестировании. \\
\hline
\texttt{f3(long double x)} & Ассемблерная функция. Возвращает значение \(f_3(x)=\frac{5}{x}\). В программе вызывается только на положительных значениях \(x\), так как по условию \(x>0\). \\
\hline
\texttt{df1(long double x)} & Ассемблерная функция. Возвращает производную \(f'_1(x)=-\frac{1.5}{(x+1)^2}\), необходимую для метода Ньютона. \\
\hline
\texttt{df2(long double x)} & Ассемблерная функция. Возвращает производную \(f'_2(x)=2.5\). \\
\hline
\texttt{df3(long double x)} & Ассемблерная функция. Возвращает производную \(f'_3(x)=-\frac{5}{x^2}\). \\
\hline
\texttt{raznost(func\_t f, func\_t g, long double x)} & Возвращает значение разности двух функций \(f(x)-g(x)\). Эта функция используется для преобразования задачи пересечения двух графиков к задаче поиска корня одной функции. \\
\hline
\texttt{raznostproisv(func\_t df, func\_t dg, long double x)} & Возвращает разность производных \(f'(x)-g'(x)\). Используется внутри метода Ньютона для функции \(F(x)=f(x)-g(x)\). \\
\hline
\texttt{root(...)} & Основная функция поиска корня уравнения \(f(x)=g(x)\) на отрезке \([a,b]\). На каждой итерации вычисляет приближение методом хорд и приближение методом Ньютона. Итерации продолжаются, пока расстояние между этими приближениями больше заданной точности \(\varepsilon\). Через указатель \texttt{iterations} возвращается число итераций. \\
\hline
\texttt{integral(func\_t f, long double a, long double b, long double eps)} & Вычисляет определенный интеграл функции \(f\) на отрезке \([a,b]\) методом трапеций. Число разбиений начинается с 4 и удваивается до тех пор, пока разность двух последовательных приближений не станет меньше \(\varepsilon\). \\
\hline
\texttt{get\_func(char *name)} & По строковому имени \texttt{"f1"}, \texttt{"f2"} или \texttt{"f3"} возвращает указатель на соответствующую функцию. Используется для обработки ключей \texttt{--test-root} и \texttt{--test-integral}. \\
\hline
\texttt{get\_dfunc(char *name)} & По имени функции возвращает указатель на ее производную: \texttt{df1}, \texttt{df2} или \texttt{df3}. Используется при тестировании функции \texttt{root}. \\
\hline
\texttt{print\_help(void)} & Выводит список доступных режимов запуска программы и формат аргументов командной строки. \\
\hline
\texttt{calculate(...)} & Выполняет основные вычисления: находит три точки пересечения, вычисляет площадь фигуры и возвращает количество итераций для каждого корня. \\
\hline
\texttt{print\_result(long double eps)} & Вызывает \texttt{calculate} для указанного значения \(\varepsilon\) и выводит корни, количество итераций и площадь. Этот режим используется при запуске программы без аргументов. \\
\hline
\texttt{print\_roots(void)} & Выводит только найденные корни для точности \(10^{-7}\). Соответствует ключу \texttt{--roots}. \\
\hline
\texttt{print\_iterations(void)} & Выводит только количество итераций для трех корней при точности \(10^{-7}\). Соответствует ключу \texttt{--iterations}. \\
\hline
\texttt{main(int argc, char *argv[])} & Главная функция программы. Анализирует аргументы командной строки, выбирает режим работы, запускает основные вычисления или тестовые режимы. При запуске без аргументов выводит результаты для \(10^{-3}\), \(10^{-5}\) и \(10^{-7}\). \\
\hline
\end{longtable}

\newpage

\section{Сборка программы (Make-файл)}

Сборка программы выполняется с помощью \texttt{make}. Ассемблерный файл \texttt{funcs.asm} транслируется NASM в 32-битный объектный файл \texttt{funcs.o}. Затем \texttt{main.c} и \texttt{funcs.o} собираются компилятором \texttt{gcc} в исполняемый файл \texttt{program}.

Схема зависимостей:
\[
\begin{array}{ccccc}
\texttt{funcs.asm} & \xrightarrow{\texttt{nasm -f elf32}} & \texttt{funcs.o} & & \\
 & & \downarrow & & \\
\texttt{main.c} & \xrightarrow{\texttt{gcc -m32}} & \texttt{program} & &
\end{array}
\]

Текст Make-файла:
\begin{verbatim}
all:
	nasm -f elf32 funcs.asm -o funcs.o
	gcc -m32 -no-pie main.c funcs.o -lm -o program

clean:
	rm -f *.o program
\end{verbatim}

Для сборки и запуска используются команды:
\begin{verbatim}
make
./program
\end{verbatim}
Для удаления объектных файлов и исполняемого файла используется команда:
\begin{verbatim}
make clean
\end{verbatim}

\newpage

\section{Отладка программы, тестирование функций}

Для проверки программы были реализованы ключи командной строки:
\begin{itemize}
    \item \texttt{--help} --- вывод справки;
    \item \texttt{--roots} --- вывод найденных корней;
    \item \texttt{--iterations} --- вывод количества итераций;
    \item \texttt{--test-root f1 f2 a b eps} --- тестирование функции \texttt{root};
    \item \texttt{--test-integral f1 a b eps} --- тестирование функции \texttt{integral}.
\end{itemize}

Проверка функции \texttt{root} проводилась на уравнениях пересечения заданных функций. Эти уравнения имеют аналитически вычисляемые корни, что позволяет сравнить численный ответ с точным.

\begin{table}[h]
\centering
\begin{tabular}{|c|c|c|c|}
\hline
Команда & Точный корень & Численный результат & Итерации \\
\hline
\texttt{--test-root f1 f2 5 6 0.001} & \(2+0.8\sqrt{15}\) & 5.098387 & 2 \\
\texttt{--test-root f1 f3 1 2 0.001} & \(\frac{1+\sqrt{241}}{12}\) & 1.377182 & 6 \\
\texttt{--test-root f2 f3 4 5 0.001} & \(\frac{19+\sqrt{561}}{10}\) & 4.268557 & 2 \\
\hline
\end{tabular}
\caption{Тестирование функции root}
\end{table}

Проверка функции \texttt{integral} проводилась на тех же функциях. Для них первообразные имеют простой аналитический вид:
\[
    \int f_1(x)\,dx=3x+1.5\ln|x+1|+C,
\]
\[
    \int f_2(x)\,dx=1.25x^2-9.5x+C,
\]
\[
    \int f_3(x)\,dx=5\ln x+C.
\]

\begin{table}[h]
\centering
\begin{tabular}{|c|c|c|}
\hline
Команда & Аналитическое значение & Численный результат \\
\hline
\texttt{--test-integral f1 1 2 0.001} & \(3+1.5\ln\frac{3}{2}\approx 3.608198\) & 3.608469 \\
\texttt{--test-integral f2 4 5 0.001} & \(1.750000\) & 1.750000 \\
\texttt{--test-integral f3 1 2 0.001} & \(5\ln2\approx 3.465736\) & 3.466041 \\
\hline
\end{tabular}
\caption{Тестирование функции integral}
\end{table}

Полученные численные значения отличаются от аналитических на величину, соответствующую заданной точности \(\varepsilon=0.001\). Также была проверена работа режимов \texttt{--help}, \texttt{--roots} и \texttt{--iterations}. Ошибок выполнения при тестировании обнаружено не было.

\newpage

\section{Программа на Си и на Ассемблере}

Исходные тексты программы имеются в архиве, приложенном к данному отчету. Архив содержит следующие файлы:
\begin{itemize}
    \item \texttt{main.c} --- основной код программы на языке Си;
    \item \texttt{funcs.asm} --- реализация функций и производных на NASM32;
    \item \texttt{Makefile} --- файл автоматической сборки.
\end{itemize}

Программа успешно собирается командой \texttt{make} и запускается командой \texttt{./program}. При запуске без аргументов программа выводит результаты вычислений для трех значений точности: \(10^{-3}\), \(10^{-5}\) и \(10^{-7}\). При запуске с ключами командной строки доступны дополнительные режимы проверки и вывода результатов.


\newpage

\section{Анализ допущенных ошибок}

В ходе выполнения задания программа была реализована в соответствии с требованиями варианта: функции и их производные вынесены в ассемблерный модуль, а основные численные алгоритмы, обработка ключей командной строки и вывод результатов реализованы на языке Си.

Во время итогового тестирования ошибок в работе программы обнаружено не было. Сборка проекта выполняется командой \texttt{make}, запуск без аргументов выводит результаты для трех значений точности, а дополнительные режимы \texttt{--roots}, \texttt{--iterations}, \texttt{--test-root} и \texttt{--test-integral} работают корректно. Полученные значения корней и интегралов согласуются с аналитическими проверками, а значение площади стабилизируется при уменьшении параметра точности \(\varepsilon\).

\newpage

\begin{raggedright}
\addcontentsline{toc}{section}{Список литературы}
\begin{thebibliography}{9}
\bibitem{stolyarov}
Столяров А.~В. Программирование на языке ассемблера NASM для ОС UNIX.

\bibitem{math}
Ильин~В.\,А., Садовничий~В.\,А., Сендов~Бл.\,Х. Математический анализ. Т.\,1~---М.: Изд-во МГУ, 1985.
\end{thebibliography}
\end{raggedright}
\end{document}
